\documentclass[11pt]{article}
\usepackage[margin=1in]{geometry}
\usepackage{amsmath,amssymb}
\usepackage{array}
\setlength{\parskip}{4pt}

\newtheorem{claim}{Claim}[section]

\title{Teaching Notes for Paper\_Draft\_1A (Sections II through V)}
\author{Companion document, not for submission}
\date{July 2026}

\begin{document}
\maketitle

\noindent A slow, pedantic walkthrough of the math-heavy sections. Each part
restates what the paper claims, then derives it with every step shown. Read
this side by side with the compiled paper. Where the paper compresses a
derivation into one line, these notes expand it into all of its lines.

\tableofcontents

% ===========================================================================
\section{Notation crib sheet}

\begin{center}
\begin{tabular}{>{$}l<{$} p{11cm}}
\mathbf{p} & a position in the plane, $\mathbf{p} = (x, y)$ \\
\mathbf{v}(\mathbf{p}) & the flow velocity at $\mathbf{p}$, components $(u, v)$ \\
\mathbf{J}(\mathbf{p}) & the $2\times2$ matrix of first derivatives of the flow \\
D(\mathbf{p}) & $\det(\mathbf{J}(\mathbf{p}))$; the single scalar landscape the whole paper navigates \\
\mathbf{g}(\mathbf{p}) & $\nabla D$, a 2-vector pointing uphill in $D$ \\
\mathbf{H}_D(\mathbf{p}) & the Hessian of $D$, symmetric $2\times2$ \\
\mathbf{f} & the flow measured at the centroid, $(u_0, v_0)$ \\
\mathbf{p}_c,\ \mathbf{p}_i & cluster centroid; position of robot $i$ \\
(u_i, v_i) & robot $i$'s flow measurement \\
\boldsymbol{\theta}_u, \boldsymbol{\theta}_v & the six fitted coefficients per component \\
\boldsymbol{\Phi} & the $6\times6$ formation matrix; row $i$ is the basis at robot $i$ \\
\rho & ring radius of the formation (pentagon circumradius) \\
(\lambda_1, \mathbf{w}_1), (\lambda_2, \mathbf{w}_2) & eigenpairs of $\mathbf{H}_D$, ordered $\lambda_1 \leq \lambda_2$ \\
\mathrm{sat}(c) & $v_{\max}\tanh(c/v_{\max})$, a smooth speed limiter \\
n, s & coordinates across ($n$) and along ($s$) the trench \\
\chi & cosine of the angle between the commanded drift and the true boundary tangent \\
\end{tabular}
\end{center}

A hat ($\hat{D}$, $\hat{\mathbf{g}}$) means the estimate built from robot
measurements; no hat means the true field quantity. The stability theorems
first assume hats equal truth; the robustness subsection puts the error back.

% ===========================================================================
\section{Problem Formulation (paper Section II)}

\subsection{Why a quadratic model}

The previous paper fit a plane to each flow component. A plane's derivatives
are constants, so a linear fit gives $\mathbf{J}$ at a point and nothing about
how $\mathbf{J}$ varies. The structures here are defined by that variation:
the Okubo-Weiss boundary is where $D = \det(\mathbf{J})$ crosses zero, and
steering to it needs $\nabla D$; the separatrix sits in a valley of $D$, and
recognizing a valley needs the second derivatives of $D$. Differentiating a
product of first derivatives produces second derivatives, so the model of
$u$ and $v$ must carry them. The cheapest such model is a quadratic:
\begin{equation}
u(\mathbf{p}_c + \tilde{\mathbf{p}}) \approx u_0 + u_x\tilde{x} + u_y\tilde{y}
+ \tfrac{1}{2}u_{xx}\tilde{x}^2 + u_{xy}\tilde{x}\tilde{y}
+ \tfrac{1}{2}u_{yy}\tilde{y}^2,
\end{equation}
and identically for $v$. The coefficients are exactly the Taylor data at the
centroid; the one-half factors exist so that this identification is exact.

\subsection{The Okubo-Weiss identity, every step}

Define, from the four entries of $\mathbf{J}$,
\begin{equation}
s_n = u_x - v_y, \qquad s_s = v_x + u_y, \qquad \omega = v_x - u_y .
\end{equation}
The Okubo-Weiss parameter is $Q = \tfrac{1}{4}(s_n^2 + s_s^2 - \omega^2)$.
Expand the three squares:
\begin{align}
s_n^2 &= u_x^2 - 2u_xv_y + v_y^2, \\
s_s^2 &= v_x^2 + 2v_xu_y + u_y^2, \\
\omega^2 &= v_x^2 - 2v_xu_y + u_y^2 .
\end{align}
Add the first two, subtract the third. The $v_x^2$ and $u_y^2$ terms cancel:
\begin{align}
s_n^2 + s_s^2 - \omega^2
&= u_x^2 - 2u_xv_y + v_y^2 + 4v_xu_y \\
&= \underbrace{(u_x + v_y)^2}_{(\nabla\cdot\mathbf{v})^2}
 - 4u_xv_y + 4u_yv_x \\
&= (\nabla\cdot\mathbf{v})^2 - 4\underbrace{(u_xv_y - u_yv_x)}_{D} .
\end{align}
(The second line adds and subtracts $2u_xv_y$ to complete the square.)
Divide by four:
\begin{equation}
Q = \tfrac{1}{4}(\nabla\cdot\mathbf{v})^2 - D .
\end{equation}
For incompressible flow $\nabla\cdot\mathbf{v} = 0$, so $Q = -D$ exactly.
The paper works with $D$, with the convention $D > 0$ inside vortex cores.

\subsection{Why the eigenvalues are $\pm\sqrt{-D}$}

The eigenvalues of a $2\times2$ matrix solve
$\lambda^2 - \operatorname{tr}(\mathbf{J})\lambda + \det(\mathbf{J}) = 0$.
Incompressibility gives $\operatorname{tr}(\mathbf{J}) = u_x + v_y = 0$, so
\begin{equation}
\lambda^2 = -D
\qquad\Longrightarrow\qquad
\lambda = \pm\sqrt{-D} \ (D<0), \qquad \lambda = \pm i\sqrt{D} \ (D>0).
\end{equation}
When $D < 0$ nearby particles separate like $e^{\sqrt{-D}\,t}$: the number
$\sqrt{-D}$ is a literal instantaneous stretching rate. This is the bridge
between the Eulerian quantity we can sense and the Lagrangian structure we
care about: more negative $D$ means stronger hyperbolicity, so the trench of
$D$ is the locus of strongest instantaneous attraction and repulsion.

\subsection{The double-gyre closed forms, derived}

The field, in field-frame coordinates $x_f = x + 1$, $y_f = y + \tfrac12$:
\begin{equation}
u = -\pi A \sin(\pi x_f)\cos(\pi y_f), \qquad
v = \pi A \cos(\pi x_f)\sin(\pi y_f) .
\end{equation}
Abbreviate $s_x = \sin\pi x_f$, $c_x = \cos\pi x_f$, and likewise for $y$.
Differentiate:
\begin{equation}
u_x = -\pi^2 A\, c_x c_y, \quad
u_y = \pi^2 A\, s_x s_y, \quad
v_x = -\pi^2 A\, s_x s_y, \quad
v_y = \pi^2 A\, c_x c_y .
\end{equation}
Note $u_x = -v_y$: the flow is divergence free, as required for $Q = -D$.
Now the determinant, keeping every sign:
\begin{align}
D &= u_x v_y - u_y v_x \\
  &= (-\pi^2 A\, c_x c_y)(\pi^2 A\, c_x c_y)
     - (\pi^2 A\, s_x s_y)(-\pi^2 A\, s_x s_y) \\
  &= -\pi^4 A^2 c_x^2 c_y^2 + \pi^4 A^2 s_x^2 s_y^2 .
\end{align}
(The old appendix dropped the minus sign on the $u_yv_x$ product and obtained
$-\pi^4A^2(c_x^2c_y^2 + s_x^2s_y^2)$, which is nonpositive everywhere. That
would say the gyre cores are not rotation dominated, which is visibly false.)

Reduce with the double-angle identities
$c^2 = \tfrac{1}{2}(1 + \cos 2\alpha)$,
$s^2 = \tfrac{1}{2}(1 - \cos 2\alpha)$.
Write $C_x = \cos 2\pi x_f$, $C_y = \cos 2\pi y_f$:
\begin{align}
c_x^2 c_y^2 &= \tfrac{1}{4}(1 + C_x)(1 + C_y)
             = \tfrac{1}{4}\bigl(1 + C_x + C_y + C_xC_y\bigr), \\
s_x^2 s_y^2 &= \tfrac{1}{4}(1 - C_x)(1 - C_y)
             = \tfrac{1}{4}\bigl(1 - C_x - C_y + C_xC_y\bigr), \\
s_x^2 s_y^2 - c_x^2 c_y^2 &= \tfrac{1}{4}\bigl(-2C_x - 2C_y\bigr)
             = -\tfrac{1}{2}(C_x + C_y) .
\end{align}
The constants and the cross terms $C_xC_y$ cancel exactly. Therefore
\begin{equation}
\boxed{\;D = -\frac{\pi^4 A^2}{2}\bigl(\cos 2\pi x_f + \cos 2\pi y_f\bigr)\;}
\label{eq:Dclosed}
\end{equation}
Differentiate (\ref{eq:Dclosed}) once and twice:
\begin{equation}
\nabla D = \pi^5 A^2
\begin{bmatrix} \sin 2\pi x_f \\ \sin 2\pi y_f \end{bmatrix},
\qquad
\mathbf{H}_D = 2\pi^6 A^2
\begin{bmatrix} \cos 2\pi x_f & 0 \\ 0 & \cos 2\pi y_f \end{bmatrix}.
\label{eq:gradHclosed}
\end{equation}
The Hessian is diagonal everywhere, i.e., $D$ is additively separable in
$x_f$ and $y_f$. Three consequences carry the whole paper.

\paragraph{(a) The zero set is a diamond.}
Apply the sum-to-product identity
$\cos P + \cos Q = 2\cos\frac{P+Q}{2}\cos\frac{P-Q}{2}$ with $P = 2\pi x_f$,
$Q = 2\pi y_f$:
\begin{equation}
D = 0
\iff \cos\pi(x_f + y_f)\,\cos\pi(x_f - y_f) = 0
\iff x_f \pm y_f \in \tfrac{1}{2} + \mathbb{Z}.
\end{equation}
Straight lines at $\pm 45^\circ$, tracing a diamond around each gyre center.
Sanity checks from (\ref{eq:Dclosed}): at a gyre center
($x_f = \tfrac12$ or $\tfrac32$, $y_f = \tfrac12$), $C_x = C_y = -1$ and
$D = +\pi^4A^2 > 0$ (rotation). At a corner saddle ($x_f = 1$, $y_f \in
\{0,1\}$), $C_x = C_y$ gives $D = -\pi^4A^2 < 0$ (strain).

\paragraph{(b) The separatrix is an exact valley of $D$.}
On the line $x_f = 1$ (world $x = 0$):
\begin{equation}
\frac{\partial D}{\partial x}\Big|_{x_f=1} = \pi^5A^2 \sin 2\pi = 0,
\qquad
\frac{\partial^2 D}{\partial x^2}\Big|_{x_f=1} = 2\pi^6A^2 \cos 2\pi
 = 2\pi^6A^2 > 0,
\end{equation}
for every $y$. Walking across the separatrix, $D$ dips to a local minimum
exactly on it: a valley line, which the paper calls a trench. Moreover the
vorticity is
\begin{equation}
\omega = v_x - u_y = -2\pi^2 A\, s_x s_y,
\end{equation}
which vanishes on $x_f = 1$ (since $s_x = \sin\pi = 0$). There $\mathbf{J}$
is symmetric, and $\sqrt{-D}$ is not merely an eigenvalue but the actual
principal strain rate. On this field the Eulerian trench and the Lagrangian
separatrix are the same curve, exactly.

\paragraph{(c) The trench floor has two wells and a crest.}
Substitute $x_f = 1$ ($C_x = 1$) into (\ref{eq:Dclosed}):
\begin{equation}
D(0, y) = -\frac{\pi^4A^2}{2}\bigl(1 + \cos 2\pi y_f\bigr).
\end{equation}
At the flow saddles ($y_f = 0$ and $1$) this is $-\pi^4A^2$: the deepest
points, the wells. At the midpoint ($y_f = \tfrac12$, the world origin) it
is exactly $0$: the crest, where the two diamonds touch. A robot descending
$D$ along the trench floor runs out of downhill at the crest, and from the
upstream side downhill points the wrong way. The ambient flow on the
separatrix runs from the upper saddle to the lower one and vanishes only at
the saddles; it is what carries the team through the crest. Keep this
picture in mind for the whole controller section.

\subsection{The problem statement}

Two tasks, both using only the six simultaneous position and velocity
samples available at each control tick: (1) get onto the trench and ride it,
with the flow, into the saddle; (2) drive $D$ to zero and circulate along
the boundary. Deliberately absent: maps, forecasts, particle integration,
memory across ticks. The estimator is memoryless.

% ===========================================================================
\section{Distributed Second-Order Estimation (paper Section III)}

\subsection{The fit and the one-half convention}

Each robot position, relative to the centroid, is pushed through the basis
\begin{equation}
\boldsymbol{\phi}(\tilde{\mathbf{p}})
= \bigl[\,1,\ \tilde{x},\ \tilde{y},\ \tfrac{1}{2}\tilde{x}^2,\
\tilde{x}\tilde{y},\ \tfrac{1}{2}\tilde{y}^2\,\bigr]^\top .
\end{equation}
Compare with the Taylor expansion of a smooth $u$ about the centroid:
\begin{equation}
u(\tilde{\mathbf{p}}) = u_0 + u_x\tilde{x} + u_y\tilde{y}
+ \tfrac{1}{2}u_{xx}\tilde{x}^2 + u_{xy}\tilde{x}\tilde{y}
+ \tfrac{1}{2}u_{yy}\tilde{y}^2 + O(\|\tilde{\mathbf{p}}\|^3).
\end{equation}
Term by term, $\boldsymbol{\phi}^\top\boldsymbol{\theta}_u$ reproduces this
with $\boldsymbol{\theta}_u = [u_0, u_x, u_y, u_{xx}, u_{xy}, u_{yy}]^\top$.
Without the halves, the fitted second-order coefficients would equal twice
the true derivatives and every downstream Hessian formula would silently
break. Stack the six basis rows into $\boldsymbol{\Phi}$ and solve
$\boldsymbol{\Phi}\boldsymbol{\theta}_u = \mathbf{u}$,
$\boldsymbol{\Phi}\boldsymbol{\theta}_v = \mathbf{v}$; one factorization
serves both since $\boldsymbol{\Phi}$ depends only on geometry.

\subsection{Six robots are necessary}

The argument is information counting made rigorous. With $N < 6$ robots,
$\boldsymbol{\Phi} \in \mathbb{R}^{N\times6}$ has a nontrivial null space:
there exists $\boldsymbol{\theta}_0 \neq \mathbf{0}$ with
$\boldsymbol{\Phi}\boldsymbol{\theta}_0 = \mathbf{0}$. Then the quadratic
fields with coefficients $\boldsymbol{\theta}_u$ and
$\boldsymbol{\theta}_u + \boldsymbol{\theta}_0$ produce byte-identical
measurements at every robot. No estimator, linear or not, can distinguish
data that are equal. Hence six robots minimum, per component, and the two
components share positions, so six robots total.

\subsection{Sufficiency is a statement about conics}

$\boldsymbol{\Phi}$ fails to be invertible exactly when some nonzero
$\mathbf{c}$ satisfies
$\boldsymbol{\phi}(\tilde{\mathbf{p}}_i)^\top\mathbf{c} = 0$ for all six
robots. But
\begin{equation}
q(\tilde{x}, \tilde{y}) = c_1 + c_2\tilde{x} + c_3\tilde{y}
+ \tfrac{c_4}{2}\tilde{x}^2 + c_5\tilde{x}\tilde{y}
+ \tfrac{c_6}{2}\tilde{y}^2
\end{equation}
is a general polynomial of degree at most two, and its zero set is a conic:
ellipse, parabola, hyperbola, or a degenerate conic such as a pair of lines
or a single line. So: $\boldsymbol{\Phi}$ singular $\iff$ all six robots lie
on one conic. A complete characterization, not merely a sufficient
condition.

Two corollaries worth internalizing.
\begin{itemize}
\item \textbf{The ring trap.} Any six robots on a single circle lie on a
conic (the circle), so a hexagon formation is singular at every size and
heading. Rings are the most natural multirobot formation and are exactly
wrong for quadratic sensing.
\item \textbf{Pentagon plus center works.} Five points in general position
(no three collinear) determine a unique conic; for the pentagon vertices
that conic is the circumscribed circle; the center robot is not on it; so no
conic passes through all six. The center robot is not a convenience, it is
the reason the matrix inverts.
\end{itemize}

\subsection{Exact noise gains, with a worked example}

If each measurement carries independent noise of standard deviation
$\sigma_{uv}$, then
$\hat{\boldsymbol{\theta}}_u - \boldsymbol{\theta}_u
= \boldsymbol{\Phi}^{-1}\boldsymbol{\eta}$, so each coefficient's standard
deviation is $\sigma_{uv}$ times the Euclidean norm of the matching row of
$\boldsymbol{\Phi}^{-1}$. For the pentagon plus center at ring radius
$\rho$ the row norms have exact values (verified symbolically):
\begin{equation}
\sigma_{u_0} = \sigma_{uv}, \qquad
\sigma_{u_x} = \sigma_{u_y} = \frac{2}{\sqrt{10}}\frac{\sigma_{uv}}{\rho},
\qquad
\sigma_{u_{xx}} = \sigma_{u_{yy}} = \frac{8}{\sqrt{10}}
\frac{\sigma_{uv}}{\rho^2},
\qquad
\sigma_{u_{xy}} = \frac{4}{\sqrt{10}}\frac{\sigma_{uv}}{\rho^2}.
\end{equation}
The value gain is $1$ because every basis function except the constant
vanishes at the origin, so interpolation forces $\hat{u}_0$ to equal the
center robot's reading verbatim. The gains are independent of formation
heading, and the gradient error covariance is isotropic,
$\tfrac{4}{10}\sigma_{uv}^2\rho^{-2}\mathbf{I}$.

Worked example at the deployed scale, $\rho = 0.075$ m and $\sigma_{uv} =
0.01$ m/s:
\begin{equation}
\sigma_{u_x} = \frac{0.6325 \times 0.01}{0.075} \approx 0.084,
\qquad
\sigma_{u_{xx}} = \frac{2.530 \times 0.01}{0.075^2} \approx 4.50 .
\end{equation}
Second derivatives are expensive: fifty times noisier than first
derivatives at this scale. This single fact is why the boundary tracker
(no Hessian needed) is intrinsically more noise tolerant than the
separatrix tracker (which lives on the Hessian eigenbasis).

\subsection{Truncation and the optimal radius}

Split the scale out of the formation matrix:
$\boldsymbol{\Phi} = \bar{\boldsymbol{\Phi}}\,\mathbf{S}_\rho$ with
$\mathbf{S}_\rho = \mathrm{diag}(1, \rho, \rho, \rho^2, \rho^2, \rho^2)$
and $\bar{\boldsymbol{\Phi}}$ shape-only. Then
\begin{equation}
\hat{\boldsymbol{\theta}} - \boldsymbol{\theta}
= \mathbf{S}_\rho^{-1}\bar{\boldsymbol{\Phi}}^{-1}\,(\text{residual vector}).
\end{equation}
For measurement noise the residual is $O(\sigma_{uv})$, giving the
$\sigma_{uv}/\rho$ and $\sigma_{uv}/\rho^2$ gains above. For model error,
if third derivatives are bounded by $M_3$, Taylor's remainder at each robot
is at most $\tfrac{M_3}{6}\rho^3$, so
\begin{equation}
|\hat{u}_x - u_x| \leq c_1 M_3 \rho^2, \qquad
|\hat{u}_{xx} - u_{xx}| \leq c_2 M_3 \rho .
\end{equation}
The two error sources pull in opposite directions. Total Hessian error:
\begin{equation}
E(\rho) = c_2 M_3 \rho + c_2' \frac{\sigma_{uv}}{\rho^2},
\qquad
E'(\rho) = c_2 M_3 - 2c_2'\frac{\sigma_{uv}}{\rho^3} = 0
\ \Longrightarrow\
\rho^* = \left(\frac{2c_2'\sigma_{uv}}{c_2 M_3}\right)^{1/3}.
\end{equation}
Read $\rho^* \propto (\sigma_{uv}/M_3)^{1/3}$ as a design rule: noisy
sensors want a bigger formation; a rapidly curving field wants a smaller
one. This is the theory behind the formation-radius sweeps, including both
observed failure modes (too small stalls in noise, too large blurs the
structure).

\subsection{The determinant perturbation identity}

For $2\times2$ matrices the determinant expansion terminates. Write
$\mathbf{J} = \begin{bmatrix} a & b \\ c & d \end{bmatrix}$ and perturb by
$\mathbf{E} = \begin{bmatrix} e_{11} & e_{12} \\ e_{21} & e_{22}
\end{bmatrix}$:
\begin{align}
\det(\mathbf{J} + \mathbf{E})
&= (a + e_{11})(d + e_{22}) - (b + e_{12})(c + e_{21}) \\
&= (ad - bc) + \bigl(a e_{22} + d e_{11} - b e_{21} - c e_{12}\bigr)
 + \bigl(e_{11}e_{22} - e_{12}e_{21}\bigr) \\
&= \det(\mathbf{J})
 + \operatorname{tr}\!\bigl(\mathrm{adj}(\mathbf{J})\,\mathbf{E}\bigr)
 + \det(\mathbf{E}),
\end{align}
where $\mathrm{adj}(\mathbf{J}) = \begin{bmatrix} d & -b \\ -c & a
\end{bmatrix}$. No higher terms exist. Cauchy-Schwarz bounds the middle
term by $\|\mathrm{adj}(\mathbf{J})\|_F\|\mathbf{E}\|_F$, and in two
dimensions $\|\mathrm{adj}(\mathbf{J})\|_F = \|\mathbf{J}\|_F$ (same four
entries up to sign and order). For the last term, from
$|e_{11}e_{22}| \leq \tfrac{1}{2}(e_{11}^2 + e_{22}^2)$ and its twin,
$|\det \mathbf{E}| \leq \tfrac{1}{2}\|\mathbf{E}\|_F^2$. Hence
\begin{equation}
|\hat{D} - D| \leq \|\mathbf{J}\|_F\|\mathbf{E}\|_F
+ \tfrac{1}{2}\|\mathbf{E}\|_F^2 :
\end{equation}
linear in the coefficient errors with slope $\|\mathbf{J}\|_F$, plus an
explicitly known quadratic remainder.

Finally, note that under the quadratic model each entry of $\mathbf{J}$ is
affine in position, so $D$ (a difference of products of affine functions)
is \emph{exactly} quadratic over the footprint. Its gradient and Hessian at
the centroid are the closed polynomial formulas in the paper; no second fit
happens.

% ===========================================================================
\section{The Controllers (paper Section IV)}

\subsection{Newton in the eigenbasis, and what is wrong with it}

On the estimated quadratic landscape the Newton step jumps to the
stationary point in one move. Using the spectral decomposition
$\mathbf{H}_D^{-1} = \sum_k \lambda_k^{-1}\mathbf{w}_k\mathbf{w}_k^\top$,
\begin{equation}
-\mathbf{H}_D^{-1}\mathbf{g}
= -\sum_{k=1}^{2}\frac{\mathbf{w}_k^\top\mathbf{g}}{\lambda_k}\,\mathbf{w}_k .
\end{equation}
Two defects. Unbounded: near-singular curvature gives arbitrarily large
commands. Indifferent: Newton solves $\nabla D = 0$ and converges to
minima, maxima, and saddles alike; near the crest the along-trench
curvature is negative and Newton deliberately climbs to the crest and
parks.

\subsection{The two repairs}

\paragraph{Saturation.} Pass each component through
$\mathrm{sat}(c) = v_{\max}\tanh(c/v_{\max})$. Since $\tanh z \approx z$
for small $z$, the exact Newton component (and its superlinear local
convergence) survives near the target; since $\tanh z \to \pm 1$, commands
flatten at $v_{\max}$ far away. One parameter sets both the cruise speed
and the width of the terminal boundary layer.

\paragraph{The eigenvalue flip.} The \emph{attract step} keeps signed
eigenvalues; the \emph{slide step} replaces $\lambda_k$ by $|\lambda_k|$:
\begin{equation}
\dot{\mathbf{p}}_c^{\,\text{des}} = \sum_k \mathrm{sat}\!\left(
-\frac{\mathbf{w}_k^\top\hat{\mathbf{g}}}{\lambda_k}\right)\mathbf{w}_k
\qquad\text{vs.}\qquad
\dot{\mathbf{p}}_c^{\,\text{des}} = \sum_k \mathrm{sat}\!\left(
-\frac{\mathbf{w}_k^\top\hat{\mathbf{g}}}{|\lambda_k|}\right)\mathbf{w}_k .
\end{equation}
In convex directions ($\lambda_k > 0$) they agree: snap to the trench
floor. In concave directions the slide step flips the sign: instead of
climbing to the stationary point it descends away from it, sliding along
the trench. This is the saddle-free Newton idea from optimization
(Dauphin et al.), repurposed as a motion primitive.

Sign ambiguity of eigenvectors is harmless: each term contains
$\mathbf{w}_k$ twice (projection and direction), so $\mathbf{w}_k \to
-\mathbf{w}_k$ changes nothing. The one place a convention is needed is the
selector, where it is imposed explicitly ($\mathbf{w}_1$ oriented downhill).

\subsection{The separatrix tracker: three laws and a selector}

The crest problem dictates the structure. Away from the separatrix, use
eigensteps. On it, $|D|$ is small and the along-trench slope vanishes at
the crest, so something else must carry the team through. That something is
the flow: the separatrix is an invariant curve, the flow runs along it and
vanishes only at the saddles. The \emph{flow step} commands the measured
flow projected on the along-trench eigenvector plus the usual Newton snap
across:
\begin{equation}
\dot{\mathbf{p}}_c^{\,\text{des}} =
\mathrm{sat}(\mathbf{f}^\top\mathbf{w}_1)\,\mathbf{w}_1
+ \mathrm{sat}\!\left(-\frac{\hat{\mathbf{g}}^\top\mathbf{w}_2}
{\lambda_2}\right)\mathbf{w}_2 .
\end{equation}
The band that activates it: $|\hat{D}| < \varepsilon_{\text{raw}}$ or
$|\hat{D}|/\|\hat{\mathbf{H}}_D\|_F < \varepsilon_{\text{dim}}$. The second
test is the important one; the ratio has units of length squared, roughly
the squared distance to the zero set under the local model, so the band
self-scales with the landscape.

Off the band, orient $\mathbf{w}_1$ downhill
($\hat{\mathbf{g}}^\top\mathbf{w}_1 \leq 0$) and ask whether the flow
agrees with descent. Yes: slide (you are downstream of the crest; descent
leads to the correct saddle). No: attract (you are upstream; the attract
step climbs to the local stationary point, which is the crest, which is
inside the band, which hands you to the flow). If $\hat{\mathbf{H}}_D$ has
no saddle structure (both eigenvalues one sign, as in the basins near the
wells), fall back to the plain attract step, which is exactly right there:
the well is a nondegenerate minimum of $D$ sitting on the flow saddle.

The one-sentence summary the proofs exploit: \emph{every mode moves the
team with the flow along the trench and applies the same Newton snap
across it.}

\subsection{The boundary tracker}

One equation, two terms, no modes. Transverse: treat $D = 0$ as scalar
root finding. The displacement $-(\hat{D}/\|\hat{\mathbf{g}}\|^2)
\hat{\mathbf{g}}$ is the smallest step that zeroes the linearized $D$
(project any candidate step onto $\hat{\mathbf{g}}$ to see minimality), and
its length $|\hat{D}|/\|\hat{\mathbf{g}}\|$ estimates the distance to the
boundary. Tangential: drift along the perpendicular of the gradient,
sign-chosen to move with the flow. With
$\hat{\mathbf{u}} = -\mathrm{sign}(\hat{D})\hat{\mathbf{g}}/
\|\hat{\mathbf{g}}\|$ and
$\hat{\mathbf{t}} = \pm\mathbf{R}_{90^\circ}\hat{\mathbf{g}}/
\|\hat{\mathbf{g}}\|$,
\begin{equation}
\dot{\mathbf{p}}_c^{\,\text{des}} =
\mathrm{sat}\!\left(\frac{|\hat{D}|}{\|\hat{\mathbf{g}}\|}\right)
\hat{\mathbf{u}}
+ \mathrm{sat}\!\left(\mathbf{f}^\top\hat{\mathbf{t}}\right)
\hat{\mathbf{t}} .
\end{equation}
Radial correction plus tangential circulation: the same skeleton as the
orbit controller of the previous paper, with boundary distance in place of
radial error.

\paragraph{Why the tangent is the gradient perpendicular.} Because that is
what a level set's tangent \emph{is}: the directional derivative of $D$
along $\hat{\mathbf{t}}$ is $\hat{\mathbf{g}}^\top\hat{\mathbf{t}} = 0$, so
motion along $\hat{\mathbf{t}}$ does not change $D$ to first order. An
earlier implementation used the $\mathbf{H}_D$ eigenvector of smallest
$|\lambda|$ instead ("the direction $D$ curves least"). Head-to-head on the
double gyre (noise-free and noisy) and on the 2 km Santa Barbara HFR field,
it lost every arena; on the double gyre its residual was four orders of
magnitude worse because on that field the two eigenvalue magnitudes
coincide exactly on the boundary (from (\ref{eq:gradHclosed}) with
$C_y = -C_x$ there), making the eigenvector ill-defined; it resolves to a
coordinate axis $45^\circ$ off the true tangent. The eigen-tangent was
removed from code and paper; this paragraph is its tombstone.

% ===========================================================================
\section{The Stability Results (paper Section V)}

\subsection{What the setup buys}

The analysis is at the navigation layer: the centroid moves exactly as
commanded, $\dot{\mathbf{p}}_c = k\,\boldsymbol{\delta}(\mathbf{p}_c)$,
with time rescaled so $k = 1$. This is a time-scale separation: the
actuator constant is $0.3$ s against transits of tens of seconds, the same
layering validated on hardware previously. Estimates are first taken exact;
the last subsection reinstates bounded error.

\subsection{Theorem 1 (attract step certificate), line by line}

\textbf{Claim.} $V_A = \tfrac{1}{2}\|\mathbf{g}(\mathbf{p}_c)\|^2$ never
increases under the attract step; trajectories converge to
$\{\mathbf{g} = \mathbf{0}\}$.

\textbf{Computation.} The gradient of $V_A$ as a function of position is
$\nabla V_A = \mathbf{H}_D\mathbf{g}$ (the Hessian is the derivative of the
gradient). Chain rule along the motion:
\begin{equation}
\dot{V}_A = \mathbf{g}^\top\mathbf{H}_D\,\dot{\mathbf{p}}_c .
\end{equation}
Substitute the attract step, written with the Newton components
$z_k = (\mathbf{w}_k^\top\mathbf{g})/\lambda_k$, so
$\dot{\mathbf{p}}_c = -\sum_k \mathrm{sat}(z_k)\mathbf{w}_k$:
\begin{align}
\dot{V}_A
&= -\sum_k \mathrm{sat}(z_k)\;\mathbf{g}^\top\mathbf{H}_D\mathbf{w}_k
 = -\sum_k \mathrm{sat}(z_k)\;\lambda_k(\mathbf{w}_k^\top\mathbf{g}) \\
&= -\sum_k \mathrm{sat}(z_k)\;\lambda_k^2 z_k
 = -v_{\max}\sum_{k}\lambda_k^2\, z_k\tanh\!\left(\frac{z_k}{v_{\max}}\right)
 \;\leq\; 0,
\end{align}
using $\mathbf{H}_D\mathbf{w}_k = \lambda_k\mathbf{w}_k$ in the second
equality and $\lambda_k(\mathbf{w}_k^\top\mathbf{g}) = \lambda_k^2 z_k$ in
the third. Each summand is a product of $\lambda_k^2 \geq 0$ and
$z\tanh z \geq 0$ ($\tanh$ has the sign of its argument). Equality forces
$z_1 = z_2 = 0$, which for nonsingular $\mathbf{H}_D$ means
$\mathbf{g} = \mathbf{0}$. LaSalle's invariance principle finishes:
bounded trajectories converge to the critical set of $D$. Near an isolated
critical point the saturations are inactive, the step is exact Newton,
$\dot{\mathbf{p}}_c \approx -(\mathbf{p}_c - \mathbf{p}^\star)$, and
convergence is locally exponential.

\textbf{What it does not say.} Only "some critical point": crest, well, or
core center. Selecting the right one is the selector's job; that is why
this theorem alone is not the whole proof.

\subsection{Lemma 1 (common transverse contraction)}

Trench coordinates: $s$ along the valley, $n$ across, normal form
\begin{equation}
D(s, n) = \varphi(s) + \tfrac{1}{2}\lambda_\perp(s)\,n^2,
\qquad \lambda_\perp \geq \underline{\lambda} > 0 .
\end{equation}
On the double gyre this is exact with $n = x$, because $D$ is additively
separable (diagonal Hessian everywhere).

The observation that does all the work: the cross-trench gradient
component is $\partial D/\partial n = \lambda_\perp n$, and the cross
eigenvalue is $\lambda_\perp > 0$. Compute the cross component of each law:
\begin{center}
\begin{tabular}{ll}
attract: & $-\dfrac{\lambda_\perp n}{\lambda_\perp} = -n$ \\[6pt]
slide: & $-\dfrac{\lambda_\perp n}{|\lambda_\perp|} = -n$
\quad (positive eigenvalue: signed and magnitude agree) \\[6pt]
flow: & the snap term was \emph{defined} as this same quantity, $-n$ .
\end{tabular}
\end{center}
Three laws, one transverse behavior:
\begin{equation}
\dot{n} = -\mathrm{sat}(n) = -v_{\max}\tanh(n/v_{\max}),
\end{equation}
which drives $n$ to zero monotonically from either side. So
$V_\perp = \tfrac{1}{2}n^2$ is a Lyapunov function \emph{common to all
modes}, and switching, however erratic, cannot break it. That is the
standard escape from dwell-time arguments for switched systems, obtained
here by construction, because the three laws were built to share their
transverse term.

\subsection{Lemma 2 (flow-consistent progress)}

\textbf{Claim.} On the trench, in every mode, the along-trench velocity
never opposes the ambient flow, and away from small balls around the
saddles its magnitude has a positive floor $m$.

Mode by mode. \emph{Flow step:} tangential term
$\mathrm{sat}(\mathbf{f}^\top\mathbf{w}_1)\mathbf{w}_1$; the projection
appears once inside sat and once as the direction, so the result is a
nonnegative multiple of the flow's tangential component. \emph{Slide:}
selected only when the flow agrees with descent, and it moves along
descent. \emph{Attract:} selected only when the flow opposes descent, and
it moves along ascent (the concave along-trench eigenvalue flips the sign
of the Newton component), which is then the flow direction. \emph{Fallback
near the wells:} descends into the well, which is where the flow is
heading, since the well sits at the stable-manifold endpoint. The floor
$m$ combines the minimum flow speed on the trench away from the saddles
with the minimum along-trench slope wherever the band is inactive; the
assumption that places the crest inside the band makes those two cases
cover the whole trench.

\subsection{Theorem 2 (traversal and capture): assembly}

Lemma 1: the tube around the trench is invariant and the team lands on the
trench exponentially, switching notwithstanding. Lemma 2: the arc-length
coordinate then advances monotonically with the flow at rate at least $m$,
so the ball around the saddle is reached in time at most
$L_\Gamma/m$. Inside the ball $\mathbf{H}_D \succ 0$ (at the double-gyre
saddles $\mathbf{H}_D = 2\pi^6A^2\,\mathbf{I}$, as nondegenerate as it
gets), the selector is in its fallback, and Theorem 1's local case gives
exponential capture.

What each assumption excludes: the normal form with
$\lambda_\perp \geq \underline{\lambda}$ excludes flat-bottomed valleys;
the eigengap excludes points where the eigenbasis itself is ill-defined;
the invariance of the trench under the flow excludes trenches that are not
flow curves (the flow step would push the team off those); the
crest-inside-band condition is the contract between the band thresholds
and the landscape.

\subsection{Theorem 3 (boundary tracking): the bracket}

With $V_1 = \tfrac{1}{2}D^2$ and the law
$\boldsymbol{\delta} = s_\perp\hat{\mathbf{u}}
+ s_{\text{tan}}\hat{\mathbf{t}}$:
\begin{equation}
\dot{V}_1 = D\,\mathbf{g}^\top\boldsymbol{\delta}
= D\,\mathbf{g}^\top\hat{\mathbf{u}}\,s_\perp
+ D\,\mathbf{g}^\top\hat{\mathbf{t}}\,s_{\text{tan}} .
\end{equation}
Perpendicular term: $\mathbf{g}^\top\hat{\mathbf{u}}
= -\mathrm{sign}(D)\|\mathbf{g}\|$ and $s_\perp \geq 0$, so it equals
$-|D|\|\mathbf{g}\|s_\perp$: always helpful. Tangential term: writing
$\chi = (\mathbf{g}/\|\mathbf{g}\|)^\top\hat{\mathbf{t}}$ and using
$|s_{\text{tan}}| \leq v_{\max}$, it is bounded by
$|D|\|\mathbf{g}\|\bar{\chi}v_{\max}$. With
$s_\perp = v_{\max}\tanh\bigl(|D|/(\|\mathbf{g}\|v_{\max})\bigr)$, factor:
\begin{equation}
\dot{V}_1 \leq -\|\mathbf{g}\|\,|D|\,v_{\max}
\left[\tanh\!\left(\frac{|D|}{\|\mathbf{g}\|v_{\max}}\right)
- \bar{\chi}\right].
\end{equation}
Everything lives in the bracket.

\paragraph{Exact estimates.} The constructed tangent satisfies
$\hat{\mathbf{t}} \perp \mathbf{g}$, so $\chi \equiv 0$ and the bracket is
positive whenever $D \neq 0$. The elementary bound
$\tanh z \geq \tanh(1)\min(z, 1)$ for $z \geq 0$ (concavity: $\tanh z / z$
is decreasing, equals $\tanh 1$ at $z = 1$) yields two regimes:
\begin{equation}
\dot{V}_1 \leq -\tanh(1)\min\bigl\{D^2,\ |D|\,\|\mathbf{g}\|v_{\max}\bigr\}.
\end{equation}
Far away ($|D| \geq \|\mathbf{g}\|v_{\max}$) the min is the second entry
and $|D|$ shrinks at fixed rate at least $\tanh(1)g_{\min}v_{\max}$: a
constant-speed reach that hits the near zone in finite time. Near the
boundary the min is $D^2$ and $\dot{V}_1 \leq -2\tanh(1)V_1$: plain
exponential decay.

\paragraph{Misaligned drift.} If errors tilt the commanded tangent by
$\bar{\chi} > 0$, the bracket changes sign at
$\tanh\bigl(|D|/(\|\mathbf{g}\|v_{\max})\bigr) = \bar{\chi}$, i.e., at
$|D| = \|\mathbf{g}\|v_{\max}\operatorname{artanh}(\bar{\chi})$. Outside
that band $V_1$ still decreases, so the band
$|D| \leq g_{\max}v_{\max}\operatorname{artanh}(\bar{\chi})$ is attractive
and invariant: practical convergence with an explicit residual.

\paragraph{Where the premise fails.} The theorem needs
$\|\mathbf{g}\| \geq g_{\min} > 0$, which fails at exactly four points per
diamond: the corners, critical points of $D$ sitting on the boundary. They
are isolated; the flow-drift fallback carries the team past them. The
theorem tells you precisely where the fallback is needed, which is why it
is not a hack.

\subsection{The ISS proposition: noise moves the target, it does not break the loop}

Every channel passes its argument through $\tanh$, which is 1-Lipschitz:
an argument error of size $d$ perturbs the command by at most $d$. So
bounded estimate errors act as a bounded additive disturbance. Transverse
channel with disturbance:
\begin{equation}
\dot{n} = -v_{\max}\tanh(n/v_{\max}) + d(t), \qquad |d| \leq \bar{d} < v_{\max}.
\end{equation}
The right side is negative whenever
$v_{\max}\tanh(|n|/v_{\max}) > \bar{d}$. Solve the equality for $|n|$:
\begin{equation}
\tanh\!\left(\frac{|n|}{v_{\max}}\right) = \frac{\bar{d}}{v_{\max}}
\iff
|n| = v_{\max}\operatorname{artanh}\!\left(\frac{\bar{d}}{v_{\max}}\right)
\approx \bar{d} \quad (\bar{d} \ll v_{\max}),
\end{equation}
using $\operatorname{artanh}(z) = z + z^3/3 + \cdots$. Above that
threshold $|n|$ shrinks; hence the ultimate bound. The boundary tracker's
version adds the $D$-estimate error into the bracket of Theorem 3,
inflating the residual band by a computable amount.

The punchline tying Sections III and V together: $\bar{d}$ is not
abstract. It splits into a noise part ($\sigma_{uv}/\rho$ or
$\sigma_{uv}/\rho^2$ via the exact gain table) and a truncation part
($M_3\rho$ or $M_3\rho^2$), so the ultimate tracking error is an explicit
function of sensor quality and formation radius, minimized at the same
$\rho^*$ as the estimation error.

\subsection{Why the boundary theorem is stronger than the trench theory}

Now stated in the paper's Discussion: the boundary is a codimension-one
regular level set, so first-order information (a gradient that does not
vanish) fully determines the transversal direction and yields a global
certificate from a single Lyapunov function. The trench lives where that
first-order information vanishes in the normal direction, so controller
and proof alike must synthesize the missing direction from curvature, and
the eigenstructure dependence, mode switching, and tube assumptions are
the price of that synthesis. Harder problem, weaker guarantees; easier
problem, stronger guarantees. The paper needs both halves to make the
point.

% ===========================================================================
\section{Questions a reviewer (or a reader) might ask}

\paragraph{Why not compute FTLE onboard?} FTLE needs particle trajectories
integrated over a horizon, which needs the field everywhere along those
trajectories, which needs a map or forecast. The team senses only its own
footprint at the current instant. The paper's position: track the Eulerian
surrogate online, validate against FTLE computed offline from the full
radar record.

\paragraph{Okubo-Weiss is not objective.} Correct, and the paper says so:
$D$ is Galilean invariant but not invariant under time-varying rotating
frames, unlike the OECS of Serra and Haller. The defense is pragmatic:
slowly evolving surface currents, one-sample computability, and the FTLE
comparison as the empirical check.

\paragraph{Why exactly six robots?} Six is proven minimal; more robots are
supported (least squares) and help with noise. Six is studied because
minimality is the scientific question and the testbed fleet is fixed.

\paragraph{What about time-varying fields?} The estimator is memoryless,
so nothing assumes a frozen field; the structures become moving targets.
The 28 hours of real radar data are the empirical answer; theory for
drifting structures is future work.

\paragraph{What if the formation deforms?} Gracefully: the fit stays exact
for quadratic fields under any deformation short of the six robots landing
on a common conic, conditioning is monitored via $\det(\boldsymbol{\Phi})$,
and the noise gains worsen continuously, not abruptly. Contrast with
fixed-weight symmetric estimators, which acquire bias proportional to the
deformation.

% ===========================================================================
\section{Where the checks live}

Every closed form in Section II and every constant in Section III was
verified symbolically (sympy) and numerically; both controllers were
verified in closed loop on the double gyre and on the 2 km Santa Barbara
HFR field. The tangent-rule comparison that removed the eigen-tangent is
recorded in the docstring of \texttt{logic\_g\_newton\_contour\_pentagon}
in
\texttt{trunk/Python\_Simulations/Vector\_Fields/VF\_Robot/src/control/pentagon\_primitives.py}
and in the git history.

\end{document}
